\documentclass[conference]{IEEEtran}
\usepackage[utf8]{inputenc}
\usepackage[T1]{fontenc}
\usepackage{cite}
\usepackage{amsmath,amssymb}
\usepackage{graphicx}
\usepackage{booktabs}
\usepackage{url}

\title{Measuring Validator Utility in Generate--Verify--Repair NetOps Pipelines: a Confirmatory Paired Study with Shared Initial Candidates}

\author{
\IEEEauthorblockN{Autonomous AI Researcher}
\IEEEauthorblockA{CNSM 2026 Agentic AI Researcher Track}
}

\begin{document}

\maketitle

\begin{abstract}
We preregistered and executed a paired confirmatory study (AC-2026-03) to evaluate whether a multidimensional validator metric suite predicts end-to-end agentic NetOps success better than a single verifier pass rate. Using a synthetic intent\ensuremath{\rightarrow}configuration generator and a hosted generate--validate--repair (GVR) adapter under shared\_initial\_candidate semantics, we ran 200 paired tasks (one shared initial generation per task reused for baseline and guarded). The deterministic validator (deterministic\_netops\_validator\_v5, v5.0) judged every sampled initial candidate valid; no repairs were applied. Baseline = 200/200, guarded = 200/200, paired difference = 0.00 (bootstrap 95\% CI [0.00, 0.00], exact McNemar two-sided p = 1.0; bootstrap\_seed = 1729, resamples = 10000). The observed ceiling invalidated the preregistered power assumptions (targeted 10 percentage-point paired improvement) and rendered the primary confirmatory test uninformative for the originally planned effect. We report preregistered design choices, precise execution accounting, representative validator-trace excerpts, quantitative analysis, failure-mode diagnostics, and artifact-grounded prescriptions for follow-up preregistered experiments (fault injection, multi-sample generation, multi-validator comparison) required to estimate validator coverage and repair value.
\end{abstract}

\section{Introduction}
Automating NetOps workflows with generative models requires reliable mechanisms to avoid harmful, silent, or wrong-but-plausible actions. A common operational pattern is generate--verify--repair (GVR): a generator (LLM) proposes a candidate configuration or plan, a deterministic validator mechanically checks correctness and policy constraints, and, when necessary, a repair module synthesizes corrected candidates or triggers human review. Prior demonstrations of GVR-like architectures in code and systems motivate their adoption in NetOps, but systematic, preregistered measurements of validator coverage (false-positive/false-negative rates), detection latency, and predictive usefulness for end-to-end guarded success are scarce [1].

We preregistered study AC-2026-03 with a paired design that isolates the effect of downstream validator-triggered repair under shared\_initial\_candidate semantics: for each task the generator produced a single initial candidate that was reused by both the baseline (no repair) and guarded (validator+repair permitted) conditions. Under these semantics any paired difference can only arise downstream from validator-triggered repair, not from independent generator sampling. The primary estimand was the paired\_success\_rate\_difference\_guarded\_minus\_baseline and the preregistered confirmatory hypothesis (H1) expected that a multidimensional validator-metric suite would explain more variance in end-to-end success than a single verifier pass rate. Our main contribution is a fully preregistered, artifact-archived experiment and an exact, transparent report of a degenerate confirmatory outcome (ceiling) with concrete, artifact-grounded prescriptions for follow-up designs that can measure validator utility.

\section{Related Work}
Work across code, systems, and NetOps communities has proposed and prototyped generate--verify--repair or closed-loop verifier architectures to reduce stochastic generator error and enable deterministic acceptance criteria [1]. Recent NetOps and AIOps benchmark efforts document substantive operational failure modes (ticket noise, false premises, and wrong-but-plausible outputs) that complicate end-to-end automation and motivate insertion of verifier layers [2],[3]. Independent system reports and preprints propose self-healing orchestrators, auditable decision traces, and verifier-guided repair to reduce silent or action-triggering failures in tool-augmented LLM systems [4]. Surveys of LLMs for telecommunications and NetOps discuss adaptation and verification trade-offs and motivate metricized evaluation of verification cost versus nominal correctness [5].

These verified records motivate a measurement focus on validator soundness and operational impact but do not provide a standardized, empirically estimable validator-metric suite tailored to NetOps GVR flows. The present preregistered experiment was designed to begin filling that measurement gap under strict preregistration and archived-execution practices; when interpreted against the cited literature we treat prior demonstrations as architectural and motivational rather than as establishing empirical coverage for NetOps validators.

\section{Methodology}
Preregistration and primary design. The study preregistration (AC-2026-03) specified a paired confirmatory design using the hosted\_netops\_gvr\_v1 adapter family with generation\_semantics = shared\_initial\_candidate and initial\_generation\_calls\_per\_task = 1. Planned task\_count = 200 pairs (planned\_episode\_count = 400), planned maximum\_model\_calls = 400, and model\_scope = gpt-5-mini via the provider bundle recorded as 'openai\_responses'. The preregistered primary estimand was paired\_success\_rate\_difference\_guarded\_minus\_baseline; the primary test was an exact two-sided McNemar with paired bootstrap percentile confidence intervals (bootstrap\_resamples = 10000, bootstrap\_seed = 1729). Targeted detectable effect was an absolute 0.10 paired improvement with \ensuremath{\alpha} = 0.05 and \textasciitilde{}80\% power under conservative planning assumptions.

Generator, tasks, and constraints. Tasks were drawn deterministically from deterministic\_netops\_task\_generator\_v5 (task\_family = intent\_configuration\_repair\_v1). Each task encodes: an initial\_state, an expected\_final\_state, an explicit intent string, allowed\_fields/allowed\_touched\_fields, workflow\_policies (e.g., minimum\_active\_in\_groups, must\_be\_down\_for\_fields), and a required\_command\_sequence. Task selection followed a preregistered deterministic index rule (challenging\_workflow\_stress\_v1). Contamination/exclusion rules (duplicate SHA-256 and embedding-similarity > 0.95) were preregistered; the execution/analysis artifacts record that no contamination exclusions occurred (analysis/contamination\_summary.csv empty).

Model and sampling configuration. The declared model version in execution/model\_configuration.json is "gpt-5-mini"; execution accounting records model\_calls\_used = 200 (one shared initial generation per pair) and maximum\_model\_calls = 400 as preregistered. Important artifact-grounded note: execution/model\_configuration.json records temperature = null (temperature = null/unset). Null/unset is explicitly not evidence of deterministic provider-side sampling and does not imply temperature = 0; provider sampling variability remains possible and is recorded in provider-call audit fields. Provider-call audit in deterministic\_reconciliation reports hosted\_model\_call\_event\_count = 200 and stage\_counts.shared\_initial\_generation = 200, consistent with the preregistered single shared generation per task.

Validator and scoring. The deterministic validator used was deterministic\_netops\_validator\_v5 (validator\_version = 5.0). The validator implements mechanistic intent-constraint checks and emits per-episode fields: normalized\_configuration, parsed\_command\_count, required\_command\_count, observed\_touched\_field\_count, violation\_count/violation\_codes, valid (boolean), repair\_applied (boolean), and repair artifacts when applicable. The preregistered binary scoring rule (full\_intent\_constraint\_validation) uses validator.valid as the success label for the confirmatory estimand. The pipeline permitted at most one repair call when validator rejected a candidate; repair events would be recorded with repair\_applied = true and repair\_provider\_trace fields.

Execution, exclusions, and missingness. The preregistered maximum\_attempts\_per\_call = 1 (no manual retries). The execution manifest reports completed\_episode\_count = 400, failed\_episode\_count = 0, and complete\_pair\_count = 200. The preregistered missingness and contamination plans were followed: confirmatory analyses used complete pairs only; artifacts include analysis/missingness\_summary.csv and analysis/contamination\_summary.csv. All raw responses, per-episode validator traces, provider-call traces, and analysis artifacts are archived in the execution bundle and referenced in this paper.

\section{Results}
Execution accounting and raw outcome summary. The run completed exactly per preregistration: planned\_episode\_count = 400, completed\_episode\_count = 400, failed\_episode\_count = 0, complete\_pair\_count = 200, and model\_calls\_used = 200 (one shared initial generation per pair). Provider-call audit entries show hosted\_model\_call\_event\_count = 200, completed\_hosted\_model\_call\_event\_count = 200, cache\_hit\_count = 0, and cache\_miss\_count = 200. The analysis condition summary (analysis/condition\_summary.csv) reports: baseline successes = 200, baseline failures = 0 (success\_rate = 1.0); guarded successes = 200, guarded failures = 0 (success\_rate = 1.0).

Primary confirmatory statistics and exact contingency. The paired contingency (analysis/paired\_contingency\_table.csv) is n\_11 = 200, n\_10 = 0, n\_01 = 0, n\_00 = 0; that is, every pair was valid under both baseline and guarded. The paired\_success\_rate\_difference\_guarded\_minus\_baseline equals 0.00. The preregistered exact McNemar two-sided test returned p = 1.0 because there are zero discordant pairs (n\_10 + n\_01 = 0). The preregistered bootstrap-percentile confidence interval (bootstrap\_resamples = 10000, bootstrap\_seed = 1729) is [0.00, 0.00], reflecting a point mass at zero in the resampled paired-difference distribution. These analysis outputs are recorded verbatim in analysis artifacts (analysis/analysis\_log.jsonl, analysis/paired\_contingency\_table.csv).

Representative validator-trace and why no repairs occurred. A typical archived validator\_trace (pair-000004) exhibits the mechanistic checks the validator enforces and illustrates why no repair was triggered:

pair\_id: pair-000004
  normalized\_configuration: "interface edge2 admin down\textbackslash{}ninterface edge2 vlan 40\textbackslash{}ninterface edge2 admin up\textbackslash{}ninterface edge1 admin down"
  parsed\_command\_count: 4
  required\_command\_count: 4
  observed\_touched\_field\_count: 3
  valid: true
  repair\_applied: false
  validator\_id: deterministic\_netops\_validator\_v5
  validator\_version: 5.0

Equivalent validator\_trace.validation\_after.valid = true and repair\_applied = false hold for every archived episode; repair\_provider\_trace fields are null for all episodes. Representative task manifests show required\_command\_count equal to parsed\_command\_count for the accepted candidates, indicating structural completeness under the validator's rules. Because baseline and guarded reused the identical shared\_initial\_candidate per pair, these per-episode validator outputs explain the degenerate paired outcome: there were no validator-detected mechanical defects to correct.

Expanded diagnostics: what the artifacts reveal about failure modes and what remains unobserved. Several archived artifacts together explain the ceiling outcome without invoking unrecorded assumptions:
  - deterministic\_reconciliation and model\_configuration.json show the single-shared-generation accounting (one initial draw per task; model\_calls\_used = 200). This sampling regime eliminated generator-sampling variance as a potential source of discordant pairs.
  - analysis/contamination\_summary.csv and analysis/missingness\_summary.csv both report no contamination flags and complete\_pairs = 200, confirming the confirmatory dataset was contamination-free and complete as preregistered.
  - validator traces uniformly report valid = true and repair\_applied = false; thus per-validator FP/FN estimates are degenerate (no false positives and no false negatives observed relative to the validator's own criteria on this sampled set), and detection-latency/repair-invocation distributions are uninformative because no repair events occurred.

Power and inferential consequences in artifact terms. The preregistered power calculation targeted detection of an absolute paired improvement of 0.10 with \ensuremath{\alpha} = 0.05 and \ensuremath{\approx}80\% power under conservative assumptions (see preregistration). Those calculations assumed non-degenerate baseline variability (for example baseline p values in \{0.3,0.4,0.5,0.6,0.7\}). The artifact-grounded observed baseline\_success\_rate = 1.0 (200/200) makes the preregistered detectable effect logically unattainable in this execution: with no observed failures in baseline there are no discordant baseline-only episodes to be corrected by the guarded condition, so the paired estimator cannot register the preregistered 10-percentage-point improvement. We therefore report the confirmatory null outcome as an artifact-valid, design-driven ceiling rather than as evidence that validators lack operational value generally.

What the archived traces allow and what they do not. The execution bundle contains full per-episode normalized\_configurations, required\_command\_count and parsed\_command\_count fields, provider-call traces, and analysis CSVs. These artifacts enable reproducible reanalyses under alternative scoring rules (for example, scoring that treats specific semantic mismatches as failures) or under changed sampling semantics (multi-sample per task) or injected faults. What the current artifacts do not support is estimation of repair effectiveness or validator false-negative rates in unperturbed production-like inputs because repair\_applied = false for every episode and no injected faults were executed in this run.

Concluding diagnostic summary. The degenerate all-valid outcome arises from the conjunction of (1) a deterministic task generator that produced candidates aligned with the validator's mechanistic constraints for the sampled seeds, (2) a single-sample-per-task shared\_initial\_candidate design that removed generator variability as a potential source of discordance, and (3) a deterministic validator that enforces structural and policy constraints but does not attempt to model operator intent beyond the encoded workflow\_policies. These artifact-grounded determinants fully account for the lack of observed repairs and the uninformative confirmatory statistic; follow-up preregistered arms (fault injection, multi-sample generation, multi-validator comparison) are required --- and are supported by the archived infrastructure --- to produce non-degenerate estimates of validator coverage, FP/FN rates, repair invocation utility, and detection latency.

\section{Discussion}
Design implications of shared\_initial\_candidate semantics. The preregistered shared\_initial\_candidate choice isolates downstream validator effects by ensuring both conditions evaluate the same initial candidate for each pair. This isolation is scientifically valuable because any observed paired difference must derive from validator-triggered repair. It also concentrates the risk of a single-sample ceiling: if the generator sample aligns with the validator for all tasks, the design yields zero observed variance in the primary estimand, as occurred here. We emphasize this property because it determines what the paired estimator can and cannot measure: it cannot detect differences attributable to generator sampling or to alternative prompt constraints when the sampling semantics were shared.

Operational interpretation and validator scope. The deterministic validator (deterministic\_netops\_validator\_v5) reliably enforces mechanistic intent-constraint checks (parsed\_command\_count, required\_command\_count, violation\_codes, valid flag). However, artifact-grounded evidence here shows that correctness-by-structure does not imply semantic intent alignment under operator expectations. That limitation is intrinsic to validators that lack an external intent model or operator feedback; detecting semantic mismatches requires either richer intent models, human-in-the-loop signals, or tailored semantic validators. The present execution does not imply validators are unnecessary---rather, it demonstrates that in a synthetic bank and single-sample regime the validator had nothing mechanical to correct.

Expanded diagnostics grounded in execution artifacts. The archived execution provides explicit numeric diagnostics that clarify why the confirmatory estimator was uninformative: model\_calls\_used = 200 (one shared generation per pair), complete\_pair\_count = 200, baseline\_success\_count = 200, guarded\_success\_count = 200, and repair\_applied = false observed for every episode. The paired contingency is a point mass (n\_11 = 200), and the bootstrap procedure (bootstrap\_seed = 1729, resamples = 10000) produced a degenerate sampling distribution with 95\% percentile CI [0.00, 0.00]. These exact values are available in the analysis artifacts (condition summary and paired contingency outputs) and show that the confirmatory null is not a statistical fluke but a deterministic consequence of the joint generator--validator--sampling configuration used.

Practical lessons for preregistered validator evaluation. From a methodological perspective, two trade-offs are clear and artifact-evident. First, shared\_initial\_candidate is a powerful control: it guarantees that any observed paired improvement is attributable to downstream validator/repair logic alone. Second, that control increases sensitivity to ceiling effects when generation produces validator-compliant candidates for the chosen seeds. For studies whose primary aim is to measure validator coverage (FP/FN) and repair utility, the preregistered design should therefore include one or both of the following preregistered modifications to avoid degeneracy: (a) a controlled fault-injection arm that deterministically introduces a known fraction f of invalid cases (syntactic, configuration-logic, or semantic-intent swaps) into the evaluation set so that the validator will face realistic defects; (b) a multi-sample generation arm with initial\_generation\_calls\_per\_task >= 3 independent draws per task to reintroduce generator variability and produce candidate diversity that can trigger validator rejections and repairs. Both prescriptions were articulated in the archived post-run prescriptions and are directly implementable within the same adapter semantics.

What validator metrics become estimable under those modifications. If faults are injected or multiple independent draws are allowed, the validator-metric suite preregistered in AC-2026-03 becomes estimable: (i) per-validator false-positive rate (FP-rate) as the fraction of validator.accept = true when the task's ground-truth is invalid; (ii) per-validator false-negative rate (FN-rate) as validator.reject = false for candidates that do actually violate the required intent sequence; (iii) mean detection latency measured from model-response receipt to validator decision (records include timestamps in provider- and validator-traces); (iv) repair-invocation rate and post-repair success fraction when repair\_applied = true is observed. These metrics map directly to the fields already emitted by the validator in the archived traces (valid, violation\_count, repair\_applied, parsed\_command\_count) and to provider-call audit counters; no additional opaque instrumentation is required to compute them once non-degenerate data are present.

Reproducibility and artifact-grounded protocol refinements. The execution artifacts document concrete accounting fields that support exact replication of the intended follow-ups: provider-call audit counts (hosted\_model\_call\_event\_count = 200), the model identifier (declared\_model\_version = gpt-5-mini), the recorded temperature = null/unset (explicitly not evidence of deterministic sampling), and the validator identifier/version (deterministic\_netops\_validator\_v5, v5.0). A follow-up preregistration should (and can, per the archived infrastructure) explicitly set: (1) initial\_generation\_calls\_per\_task = k (e.g., k = 3) with independent generation seeds recorded for each draw; (2) a deterministic fault-injection schedule specifying fault classes and the exact fraction f of tasks affected; and (3) a plan to measure detection latency using the existing timestamps. These concrete protocol elements are already supported by the adapter and manifest fields and would remove ambiguity about sampling and allow FP/FN estimation without inventing new instrumentation.

Threats to validity revisited with artifact evidence. Internal validity: the shared\_initial\_candidate design intentionally removes generator sampling as a confound but introduces a ceiling risk; the exact accounting shows that risk materialized (complete agreement across conditions). Construct validity: validator.valid captures mechanistic constraint satisfaction but not semantic intent alignment---this remains the central construct gap and explains why correct-by-structure outputs can still be operationally unacceptable. External validity: the study executed one model (gpt-5-mini) and one deterministic validator against a synthetic task bank; therefore caution is required when extrapolating to multi-vendor production networks or to models and validators with different coverage. Conclusion validity: the archived bootstrap and contingency outputs (seed/resamples specified) demonstrate the degenerate inferential state (zero variance) and justify treating the confirmatory hypothesis as untestable under these run parameters rather than as a failed hypothesis test.

Operational implications and recommended reporting practices. For experimenters and operators evaluating GVR pipelines, we recommend preregistering both (a) the sampling semantics (shared\_initial\_candidate vs independent draws) and (b) explicit anti-ceiling design elements (fault injection or multi-draw sampling) whenever the evaluation goal is to estimate validator coverage or repair utility. Empirically reporting the exact counts we used (model\_calls\_used, complete\_pair\_count, baseline/guarded success counts, repair\_applied counts, bootstrap\_seed/resamples) increases interpretability because readers can immediately diagnose whether a noninformative null arose from lack of validator failures or from insufficient sample diversity.

Summary. The execution artifacts and numeric diagnostics make the outcome clear: under the chosen shared\_initial\_candidate semantics, the combination of deterministic seeds, single-sample generation, and a mechanistic validator that checks structural constraints produced a ceiling effect (200/200 valid). That outcome is scientifically informative about the joint behavior of generator+validator under these settings and usefully constrains what conclusions can and cannot be drawn. It also supplies a concrete, artifact-grounded roadmap---fault injection, multi-sample generation, and multi-validator comparison---that will enable future preregistered studies to estimate the validator metrics of interest and to quantify repair utility in operationally meaningful conditions.

\section{Limitations}
\begin{itemize}
\item Shared\_initial\_candidate semantics constrained inference: baseline and guarded reused the same initial candidate per pair; any paired difference could only arise from validator-triggered repair (not from independent generator sampling).
\item Synthetic task bank (difficulty 'very\_hard'/'extreme') is not equivalent to heterogeneous production networks (multi-vendor devices, real ticket noise); external validity is limited.
\item Single-model experiment: only gpt-5-mini was executed; no multi-model generality claims are made.
\item No repairs were applied because all initial candidates passed the deterministic validator; therefore the study could not estimate repair effectiveness or validator false-negative behavior in this execution.
\item Validator scope: deterministic\_netops\_validator\_v5 enforces mechanistic intent-constraint checks but does not guarantee detection of semantic intent misalignment (correct-by-structure but incorrect-by-intent).
\item Recorded execution/model\_configuration.json temperature = null/unset does not imply deterministic sampling; provider-side sampling behavior may vary and is documented in execution manifests.
\end{itemize}

\section{Conclusion}
We executed a fully preregistered paired confirmatory study (AC-2026-03) under shared\_initial\_candidate semantics to measure validator utility in NetOps GVR pipelines. For the synthetic task bank, the sampled gpt-5-mini generations, and deterministic\_netops\_validator\_v5 used here, every sampled initial candidate passed the validator's mechanistic checks so no repairs were applied and guarded success equaled baseline (200/200). The preregistered power assumptions (targeted 10-percentage-point improvement) were invalidated by this ceiling outcome. The result is artifact-grounded and reproducible from the archived bundle. We publish the exact execution and analysis artifacts to support replication and provide concrete, preregistered experiment prescriptions (fault injection, multi-sample generation, multi-validator comparison) that are required to produce estimable validator FP/FN behavior and repair value in operationally meaningful conditions.

\section*{Disclosure Statement}
Disclosure (concise): Declared model and provider: gpt-5-mini via provider bundle 'openai\_responses'. Orchestration/adapter: hosted\_netops\_gvr\_v1. Deterministic validator: deterministic\_netops\_validator\_v5 (validator\_version = 5.0). Preregistration: AC-2026-03 (preregistration\_sha256 = de37b485\allowbreak{}421cb089\allowbreak{}5a98df38\allowbreak{}b6b3baad\allowbreak{}1dc7c13c\allowbreak{}81065e3d\allowbreak{}8240d2a9\allowbreak{}4a0ff844). Master prompt immutable reference: provenance/master\_prompt.txt, SHA-256 = 1872df1e\allowbreak{}1805d2d9\allowbreak{}6940456c\allowbreak{}a016bd66\allowbreak{}5d1d5196\allowbreak{}add77f5a\allowbreak{}cdf1582b\allowbreak{}b39b15ba. Autonomy boundary: a human Research Director selected the broad topic family and approved the master prompt before the autonomy lock; after the lock the AI system autonomously performed literature review, preregistration, experiment design, execution, analysis, internal peer review and manuscript drafting without manual editing of technical content. Sampling/generation semantics: shared\_initial\_candidate (one initial sampled generation per task reused for baseline and guarded); execution/model\_configuration.json recorded temperature = null/unset (null/unset is not evidence of deterministic sampling and does not imply temperature = 0). Call budget and usage (preregistered): maximum\_model\_calls = 400; actual model\_calls\_used = 200 (one shared initial generation per pair). All raw outputs, per-episode validator traces, execution manifests, and analysis artifacts are archived in the supplied execution bundle (see execution/execution\_manifest.json). No human manually edited the manuscript technical content after the autonomy lock.

\begin{thebibliography}{99}
\bibitem{ref1} [1] R. Cadenas, "Convergent Correctness in Stochastic Code Generation: A Generate-Verify-Repair Architecture for Deterministic Validation of Autonomous AI Coding Agents in Regulated Environments," SSRN Electronic Journal, 2026. doi:10.2139/ssrn.6754899.
\bibitem{ref2} [2] K.-H. Tseng, N. Bogahawatta, Y. Ginige, K. Patel, K. Dakic, S. Seneviratne, "FaulT-Bench: Towards Benchmarking Network Troubleshooting LLM Agents under Unreliable User Tickets," arXiv:2608.27021, 2026.
\bibitem{ref3} [3] Y. Liu, C. Pei, L. Xu, B. Chen, M. Sun, Z. Zhang, Y. Sun, S. Zhang, K. Wang, H. Zhang, J. Li, G. Xie, X. Wen, X. Nie, M. Ma, P. Dan, "OpsEval: A Comprehensive IT Operations Benchmark Suite for Large Language Models," arXiv:2310.07637, 2023.
\bibitem{ref4} [4] R. S. Babu and A. Agrawal, "Self-Healing Agentic Orchestrators for Reliable Tool-Augmented Large Language Model Systems," arXiv:2606.01416, 2026.
\bibitem{ref5} [5] H. Zhou, C. Hu, Y. Yuan, Y. Cui, Y. Jin, C. Chen, H. Wu, D. Yuan, L. Jiang, D. Wu, X. Liu, J. Zhang, X. Wang, J. Liu, "Large Language Model (LLM) for Telecommunications: A Comprehensive Survey on Principles, Key Techniques, and Opportunities," IEEE Commun. Surveys \& Tutorials, 2024. doi:10.1109/COMST.2024.3465447.
\end{thebibliography}

\end{document}
